% Intermediate summary matrix for Q6. May later be condensed, moved to an appendix, or redrawn as a figure.
\begin{table}[tbp]
\centering
\footnotesize
\singlespacing
\caption{Working summary for Q6: habituation and recovery. Four processes can yield similar declining curves; only some have diagnostic tests in the current literature.}
\label{tab:summary-q6}
\begin{tabular}{@{}p{1.5in}p{3.7in}p{1.0in}@{}}
\toprule
\textbf{Claim or process} & \textbf{Evidence} & \textbf{Status} \\
\midrule
Stimulus-/caller-specific decline & Field sciurids: decline to one caller; recovery when caller changes & Strong \\
\addlinespace[2pt]
Significance slows habituation rate & Needs natural vs matched synthetic under identical schedules & Thin / gap \\
\addlinespace[2pt]
Significance speeds recovery & More recovery to loud natural distress squeal than mimic after rest & Thin (1 study) \\
\addlinespace[2pt]
Significance prevents habituation & Prosocial 50-kHz approach habituates even when playback predicts interaction & Strong against \\
\addlinespace[2pt]
Learned unreliability & Caller history without predators (see Q5); can mimic habituation curves & Confound \\
\addlinespace[2pt]
Motivational decline & Possible when outcomes are absent or devalued; rarely isolated & Underspecified \\
\addlinespace[2pt]
Temporal expectation & Schedule effects (see Q7); must hold exposure constant & Separate account \\
\bottomrule
\end{tabular}
\end{table}
